\section{Variational Inference vs.\ EM}

Variational Inference (VI) and the Expectation--Maximization (EM) algorithm are closely related. Both arise from rewriting the log-likelihood using an auxiliary distribution over latent variables, and both can be interpreted as optimization procedures on a lower bound. However, they are not identical: EM is a special and more restrictive procedure, whereas variational inference is a broader framework.

\subsection{Latent-variable setting}

Suppose \(x\) denotes observed data, \(z\) latent variables, and \(\theta\) model parameters. The joint model is
\[
p(x,z \mid \theta),
\]
and the observed-data likelihood is
\[
p(x\mid \theta)=\sum_z p(x,z\mid \theta)
\]
in the discrete case, or
\[
p(x\mid \theta)=\int p(x,z\mid \theta)\,dz
\]
in the continuous case.

The difficulty is that direct maximization of
\[
\log p(x\mid \theta)
\]
is often hard because the latent variable \(z\) is unobserved.

\subsection{A common variational identity}

Introduce any distribution \(q(z)\) over the latent variable. Then
\[
\log p(x\mid \theta)
=
\mathcal{L}(q,\theta)
+
\mathrm{KL}\!\left(q(z)\,\|\,p(z\mid x,\theta)\right),
\]
where
\[
\mathcal{L}(q,\theta)
=
\mathbb{E}_{q(z)}[\log p(x,z\mid \theta)]
-
\mathbb{E}_{q(z)}[\log q(z)].
\]

Since Kullback--Leibler divergence is nonnegative,
\[
\mathrm{KL}\!\left(q(z)\,\|\,p(z\mid x,\theta)\right)\ge 0,
\]
we obtain the lower bound
\[
\mathcal{L}(q,\theta)\le \log p(x\mid \theta).
\]
This lower bound is often called the \emph{evidence lower bound} (ELBO).

Thus both EM and VI are based on the same idea:
\begin{quote}
replace the difficult objective \(\log p(x\mid \theta)\) by a tractable lower bound \(\mathcal{L}(q,\theta)\).
\end{quote}

\subsection{Expectation--Maximization}

EM alternates between two steps.

\subsubsection*{E-step}
Given the current parameter value \(\theta^{(t)}\), set
\[
q^{(t+1)}(z)=p(z\mid x,\theta^{(t)}).
\]
This is the exact posterior over the latent variables under the current model.

Because of this exact choice,
\[
\mathrm{KL}\!\left(q^{(t+1)}(z)\,\|\,p(z\mid x,\theta^{(t)})\right)=0,
\]
so the lower bound is tight at \(\theta^{(t)}\).

\subsubsection*{M-step}
Then maximize the bound with respect to \(\theta\):
\[
\theta^{(t+1)}
=
\arg\max_{\theta}\mathcal{L}(q^{(t+1)},\theta).
\]
Since the entropy term in \(\mathcal{L}(q,\theta)\) does not depend on \(\theta\), this is equivalent to maximizing
\[
Q(\theta\mid \theta^{(t)})
=
\mathbb{E}_{p(z\mid x,\theta^{(t)})}
\big[\log p(x,z\mid \theta)\big].
\]

Hence EM can be written as:
\[
\text{E-step: exact posterior} \qquad
\text{M-step: maximize expected complete-data log-likelihood.}
\]

\subsection{Variational Inference}

Variational inference uses the same lower bound
\[
\mathcal{L}(q,\theta)
=
\mathbb{E}_{q(z)}[\log p(x,z\mid \theta)]
-
\mathbb{E}_{q(z)}[\log q(z)],
\]
but does \emph{not} require \(q(z)\) to equal the exact posterior.

Instead, one chooses a tractable family of distributions
\[
\mathcal{Q}=\{q_\phi(z)\},
\]
often parametrized by variational parameters \(\phi\), and solves
\[
q_\phi^*(z)
=
\arg\max_{q_\phi\in \mathcal Q}\mathcal{L}(q_\phi,\theta),
\]
or equivalently
\[
q_\phi^*(z)
=
\arg\min_{q_\phi\in \mathcal Q}
\mathrm{KL}\!\left(q_\phi(z)\,\|\,p(z\mid x,\theta)\right).
\]

Thus VI performs approximate posterior inference:
\begin{quote}
instead of using the exact posterior, it finds the best approximation within a restricted family.
\end{quote}

\subsection{Core difference}

The essential difference is:

\begin{itemize}
    \item \textbf{EM:} the E-step uses the \emph{exact} posterior
    \[
    q(z)=p(z\mid x,\theta).
    \]
    \item \textbf{VI:} one restricts \(q\) to a tractable family \(\mathcal Q\), so generally
    \[
    q(z)\neq p(z\mid x,\theta).
    \]
\end{itemize}

Therefore:

\begin{quote}
EM is exact posterior updating plus parameter maximization, whereas variational inference is approximate posterior optimization over a chosen family.
\end{quote}

\subsection{EM as a special case of variational inference}

EM may be viewed as a special case of variational inference in which the variational family is unrestricted and rich enough to contain the true posterior exactly. In that case, the optimal \(q\) is
\[
q^*(z)=p(z\mid x,\theta),
\]
and the KL divergence becomes zero in every E-step.

Thus:
\[
\boxed{\text{EM is a special case of variational optimization.}}
\]

\subsection{Why VI is broader}

VI becomes necessary when the exact posterior
\[
p(z\mid x,\theta)
\]
is itself intractable. This happens in many modern latent-variable models, especially in Bayesian models and deep generative models.

In such settings, EM cannot be carried out exactly because the E-step is impossible in closed form. Variational inference then replaces the exact E-step by an approximate optimization:
\[
q_\phi(z)\approx p(z\mid x,\theta).
\]

Hence VI extends the EM philosophy to settings where exact posterior computation is no longer feasible.

\subsection{Optimization viewpoint}

Both methods optimize the same lower bound, but differently:

\begin{itemize}
    \item \textbf{EM:} coordinate ascent on \(\mathcal{L}(q,\theta)\) with an exact update in \(q\).
    \item \textbf{VI:} coordinate ascent or gradient-based optimization on \(\mathcal{L}(q_\phi,\theta)\) within a restricted family.
\end{itemize}

In EM, the bound touches the true log-likelihood at each E-step. In VI, the bound is generally not tight because the approximate family cannot exactly represent the posterior.

\subsection{Comparison table}

\[
\begin{array}{|l|l|l|}
\hline
\textbf{Aspect} & \textbf{EM} & \textbf{Variational Inference} \\
\hline
Auxiliary object & q(z) & q_\phi(z) \\
\hline
Posterior treatment & \text{exact} & \text{approximate} \\
\hline
E-step / inference step &
q(z)=p(z\mid x,\theta) &
q_\phi(z)\approx p(z\mid x,\theta) \\
\hline
Optimization target & \mathcal{L}(q,\theta) & \mathcal{L}(q_\phi,\theta) \\
\hline
KL term after inference update & 0 & \text{usually } > 0 \\
\hline
Applicability & \text{when exact posterior is tractable} &
\text{when exact posterior is intractable} \\
\hline
Interpretation & \text{exact latent-variable ML} &
\text{approximate inference / bound optimization} \\
\hline
\end{array}
\]

\subsection{Conceptual summary}

The best concise relationship is:
\begin{quote}
EM and VI are built from the same variational lower-bound identity. EM uses the exact posterior in the auxiliary step, while variational inference replaces it by the best approximation in a restricted family.
\end{quote}

Equivalently:
\[
\boxed{
\text{EM} = \text{exact variational coordinate ascent},
\qquad
\text{VI} = \text{approximate variational optimization}.
}
\]

\subsection{Final remark}

In modern machine learning, VI is often discussed mainly as an approximate Bayesian inference method, while EM is often discussed mainly as a maximum-likelihood algorithm for latent-variable models. Mathematically, however, they are very closely connected: both optimize a lower bound obtained by introducing a distribution over latent variables.
